
\documentclass{article}
\usepackage{colortbl}
\usepackage{makecell}
\usepackage{multirow}
\usepackage{supertabular}

\begin{document}

\newcounter{utterance}

\centering \large Interaction Transcript for game `cladder', experiment `full\_v1.5\_default', episode 4591 with llama{-}3.1{-}8b{-}instruct.
\vspace{24pt}

{ \footnotesize  \setcounter{utterance}{1}
\setlength{\tabcolsep}{0pt}
\begin{supertabular}{c@{$\;$}|p{.15\linewidth}@{}p{.15\linewidth}p{.15\linewidth}p{.15\linewidth}p{.15\linewidth}p{.15\linewidth}}
    \# & \multicolumn{2}{c}{Player} && \multicolumn{2}{c}{Game Master} \\
    \hline

    \theutterance \stepcounter{utterance}  
    & & & \multicolumn{4}{p{0.6\linewidth}}{
        \cellcolor[rgb]{0.9,0.9,0.9}{
            \makecell[{{p{\linewidth}}}]{
                \texttt{\tiny{[P1$\langle$GM]}}
                \texttt{You are an expert in causal inference. The following question is not a typical commonsense query, but rather a meticulously designed question created by a professor specializing in causal inference, intended to assess the students' mastery of the course content.} \\
\texttt{Imagine a self{-}contained, hypothetical world with only the following conditions, and without any unmentioned factors or causal relationships: Having a brother has a direct effect on wife and alarm clock. Wife has a direct effect on alarm clock. The overall probability of having a brother is 1\%. The probability of not having a brother and ringing alarm is 64\%. The probability of having a brother and ringing alarm is 0\%. Is the chance of ringing alarm larger when observing having a brother?} \\
\texttt{Start your answer with 'Yes' or 'No', followed by additional reasoning or evidence to support your explanation.} \\
            }
        }
    }
     \\ \\

    \theutterance \stepcounter{utterance}  
    & \multicolumn{4}{p{0.6\linewidth}}{
        \cellcolor[rgb]{0.561,0.737,0.561}{
            \makecell[{{p{\linewidth}}}]{
                \texttt{\tiny{[P1$\rangle$GM]}}
                \texttt{No.} \\
\\ 
\texttt{To answer this question, we need to apply the principles of causal inference, specifically the concept of causal effects and the potential outcomes framework. } \\
\\ 
\texttt{Let's denote the presence of a brother as B, the presence of a ringing alarm as A, and the presence of a wife as W. We are given the following information:} \\
\\ 
\texttt{1. P(B) = 0.01 (probability of having a brother)} \\
\texttt{2. P(¬B ∧ A) = 0.64 (probability of not having a brother and ringing alarm)} \\
\texttt{3. P(B ∧ A) = 0 (probability of having a brother and ringing alarm)} \\
\\ 
\texttt{We want to determine if the presence of a brother (B) increases the probability of a ringing alarm (A). To do this, we need to compare the probability of a ringing alarm when we observe having a brother (P(A|B)) to the probability of a ringing alarm when we do not observe having a brother (P(A|¬B)).} \\
\\ 
\texttt{Using the law of total probability, we can express P(A|B) as:} \\
\\ 
\texttt{P(A|B) = P(A|B)P(B) + P(A|¬B)P(¬B)} \\
\\ 
\texttt{Since P(B ∧ A) = 0, we know that P(A|B) = 0 (because if having a brother implies a ringing alarm, then P(B ∧ A) would be greater than 0). Therefore:} \\
\\ 
\texttt{P(A|B) = 0} \\
\\ 
\texttt{Now, let's calculate P(A|¬B):} \\
\\ 
\texttt{P(A|¬B) = P(A|¬B)P(¬B) + P(A|B)P(B)} \\
\texttt{= P(A|¬B)P(¬B) + 0} \\
\texttt{= P(A|¬B)P(¬B)} \\
\\ 
\texttt{We are given that P(¬B ∧ A) = 0.64, which implies that P(A|¬B)P(¬B) = 0.64. Since P(¬B) = 0.99, we can solve for P(A|¬B):} \\
\\ 
\texttt{P(A|¬B) = 0.64 / 0.99} \\
\texttt{= 0.6465} \\
\\ 
\texttt{Comparing P(A|B) = 0 to P(A|¬B) = 0.6465, we see that the probability of a ringing alarm is actually larger when we do not observe having a brother. Therefore, the presence of a brother does not increase the probability of a ringing alarm.} \\
            }
        }
    }
    & & \\ \\

    \theutterance \stepcounter{utterance}  
    & & & \multicolumn{2}{p{0.3\linewidth}}{
        \cellcolor[rgb]{0.9,0.9,0.9}{
            \makecell[{{p{\linewidth}}}]{
                \texttt{\tiny{[GM$|$GM]}}
                \texttt{No} \\
            }
        }
    }
    & & \\ \\

    \theutterance \stepcounter{utterance}  
    & & & \multicolumn{2}{p{0.3\linewidth}}{
        \cellcolor[rgb]{0.9,0.9,0.9}{
            \makecell[{{p{\linewidth}}}]{
                \texttt{\tiny{[GM$|$GM]}}
                \texttt{no} \\
            }
        }
    }
    & & \\ \\

    \theutterance \stepcounter{utterance}  
    & & & \multicolumn{2}{p{0.3\linewidth}}{
        \cellcolor[rgb]{0.9,0.9,0.9}{
            \makecell[{{p{\linewidth}}}]{
                \texttt{\tiny{[GM$|$GM]}}
                \texttt{game\_result = WIN} \\
            }
        }
    }
    & & \\ \\

\end{supertabular}
}

\end{document}
